\chapter{Solutions to the Schrödinger's equation}

This chapter explores general solutions to the Schrödinger equation and their physical interpretation.

\section{Stationary solution}

The separation of variables technique is a powerful method for solving partial differential equations like the Schrödinger equation. This approach assumes that solutions can be expressed as products of functions, each depending on a single variable.


The separation of variables method assumes a solution of the form:

\begin{equation}
\Psi(\vec{x}, t) = f(t)\Psi(\vec{x})
\end{equation}

Substituting this into the Schrödinger equation:

\begin{equation}
i\hbar\Psi\frac{\partial f}{\partial t} = -\frac{\hbar^2}{2m}f\nabla^2\Psi + Uf\Psi
\end{equation}

Dividing by $f\Psi$ (assuming both are non-zero):

\begin{align}
\frac{1}{f(t)}i\hbar\frac{\partial f}{\partial t} &= \frac{1}{\Psi}\left(-\frac{\hbar^2}{2m}\nabla^2\Psi + U\Psi\right) \\
\frac{1}{f(t)}i\hbar\frac{\partial f}{\partial t} &= \frac{1}{\Psi}\hat{H}\Psi
\end{align}

Since the left side depends only on time and the right side only on spatial coordinates, both must equal a constant, which we denote as $E$. This leads to two separate equations:

\[
\left\{\begin{array}{l}
i\hbar\frac{\partial f}{\partial t} = Ef(t) \\
\hat{H}\Psi = E\Psi
\end{array}\right. \Rightarrow \left\{\begin{array}{l}
f(t) = e^{-\frac{iE}{\hbar}t} \\
\hat{H}\Psi_E = E\Psi_E
\end{array}\right.
\]

The first equation has the solution $f(t) = e^{-\frac{iE}{\hbar}t}$, while the second is an eigenvalue equation for the Hamiltonian operator. The energy $E$ represents an eigenvalue, and $\Psi_E$ the corresponding eigenfunction.

These solutions, called stationary states, have special significance in quantum mechanics as they represent states of definite energy. While the spatial probability distribution $|\Psi_E(\vec{x})|^2$ remains constant in time, the phase evolves uniformly.

\section{Superposition principle}

The stationary solution has the form:

\begin{equation}
\Psi(\vec{x}, t) = \exp\left(-\frac{iE}{\hbar}t\right)\Psi_E(\vec{x})
\end{equation}

Since the Schrödinger equation is linear, any linear combination of solutions is also a solution. This leads to the superposition principle, allowing us to write a general solution as:

\begin{equation}
\Psi(\vec{x}, t) = \sum_n \underbrace{c(E_n)}_{c_n}\exp\left(-\frac{iE_n}{\hbar}t\right)\Psi_{E_n}(\vec{x})
\end{equation}

Where $c_n$ are complex coefficients determining the contribution of each eigenstate to the total wavefunction.

To verify this is indeed a solution, we can substitute it into the Schrödinger equation. For the time derivative:


Verifying the superposition principle by substituting into the Schrödinger equation:

For the time derivative:
\begin{align}
-i\hbar\frac{\partial\Psi(\vec{x},t)}{\partial t} &= -i\hbar\frac{\partial}{\partial t}\sum_n c_n\exp\left(-\frac{iE_n}{\hbar}t\right)\Psi_{E_n}(\vec{x}) \\
&= \sum_n c_n\Psi_{E_n}(\vec{x})(-i\hbar)\frac{\partial}{\partial t}\left[\exp\left(-\frac{iE_n}{\hbar}t\right)\right] \label{eq:2.7} \\
&= \sum_n c_nE_n\exp\left(-\frac{iE_n}{\hbar}t\right)\Psi_{E_n}(\vec{x})
\end{align}

For the Hamiltonian operating on the wavefunction:
\begin{equation}
\hat{H}\Psi(\vec{x},t) = \hat{H}\sum_n c_n\exp\left(-\frac{iE_n}{\hbar}t\right)\Psi_{E_n}(\vec{x})
\end{equation}

Since the Hamiltonian operates only on the spatial part:
\begin{equation}
\hat{H}\Psi(\vec{x},t) = \sum_n c_n\exp\left(-\frac{iE_n}{\hbar}t\right)\hat{H}\Psi_{E_n}(\vec{x})
\end{equation}

Using the eigenvalue equation $\hat{H}\Psi_{E_n} = E_n\Psi_{E_n}$:
\begin{equation}
\hat{H}\Psi(\vec{x},t) = \sum_n c_n\exp\left(-\frac{iE_n}{\hbar}t\right)E_n\Psi_{E_n}(\vec{x}) \label{eq:2.10}
\end{equation}

Comparing equations \eqref{eq:2.7} and \eqref{eq:2.10}, we see they are identical, confirming that the superposition indeed satisfies the Schrödinger equation.

\section{Formal solution to the Schrödinger's equation}

The time evolution of quantum states can be elegantly expressed using the quantum propagator, which connects the initial state to its future evolution:

\begin{equation}
\underbrace{\Psi(\vec{x},t)}_{\text{propagated state}} = \underbrace{\exp\left(-\frac{i}{\hbar}\hat{H}t\right)}_{\text{propagator}}\underbrace{\Psi(\vec{x},0)}_{\text{initial state}}
\end{equation}

The exponential operator can be expanded as a Taylor series:

\begin{align}
\Psi(\vec{x},t) = \exp\left(-\frac{i}{\hbar}\hat{H}t\right)\Psi(\vec{x},0) &\stackrel{\text{Taylor}}{=} \sum_{n=0}^{\infty}\frac{\left(-\frac{i}{\hbar}t\right)^n}{n!}\hat{H}^n\Psi(\vec{x},0) \\
&= \sum_{n=0}^{\infty}\frac{\left(-\frac{i}{\hbar}t\right)^n}{n!}E^n\Psi(\vec{x},0)  \\
&= \exp\left(-\frac{i}{\hbar}Et\right)\Psi(\vec{x},0)
\end{align}

This formal solution demonstrates how the time-evolution operator $\exp\left(-\frac{i}{\hbar}\hat{H}t\right)$ acts as a propagator that evolves the initial state forward in time. For stationary states (eigenstates of the Hamiltonian), this simply adds a time-dependent phase factor.

For a general initial state that is a superposition of energy eigenstates, the propagator acts on each component independently, consistent with our earlier derivation using the superposition principle.

The propagator formalism offers several advantages:

1. It provides a compact mathematical expression for time evolution
2. It preserves probability (unitarity) by construction
3. It can be extended to more complex scenarios, including time-dependent Hamiltonians
4. It connects quantum mechanics to path integral formulations

This approach also highlights the fundamental role of the Hamiltonian as the generator of time translations in quantum mechanics, a direct consequence of energy conservation and time-translation symmetry.


The formal solution to the Schrödinger equation can be verified by direct substitution. For a superposition of energy eigenstates evolved by the propagator:

\begin{align}
i\hbar\frac{\partial}{\partial t}\sum_n c_n\exp\left(-\frac{i}{\hbar}\hat{H}t\right)\Psi_{E_n} &= \sum_n c_n i\hbar\frac{\partial}{\partial t}\exp\left(-\frac{i}{\hbar}\hat{H}t\right)\Psi_{E_n} \\
&= \sum_n c_n\exp\left(-\frac{i}{\hbar}\hat{H}t\right)\hat{H}\Psi_{E_n} \\
&= \sum_n c_n\exp\left(-\frac{i}{\hbar}\hat{H}t\right)E_n\Psi_{E_n} \label{eq:2.13} \\
&= \exp\left(-\frac{i}{\hbar}\hat{H}t\right)\sum_n c_nE_n\Psi_{E_n}
\end{align}

And for the Hamiltonian acting on the time-evolved state:
\begin{align}
\hat{H}\sum_n c_n\exp\left(-\frac{i}{\hbar}\hat{H}t\right)\Psi_{E_n} &= \sum_n c_n\hat{H}\exp\left(-\frac{i}{\hbar}\hat{H}t\right)\Psi_{E_n} \\
&= \sum_n c_n\exp\left(-\frac{i}{\hbar}\hat{H}t\right)\hat{H}\Psi_{E_n} \\
&= \sum_n c_n\exp\left(-\frac{i}{\hbar}\hat{H}t\right)E_n\Psi_{E_n} \label{eq:2.14} \\
&= \exp\left(-\frac{i}{\hbar}\hat{H}t\right)\sum_n c_nE_n\Psi_{E_n}
\end{align}

Since the right-hand sides of equations \eqref{eq:2.13} and \eqref{eq:2.14} are identical, the propagator formulation satisfies the Schrödinger equation.

Alternatively, we can verify this using the Taylor expansion of the propagator:
\begin{align}
i\hbar\frac{\partial}{\partial t}\sum_{n=0}^{\infty}\frac{\left(-\frac{i}{\hbar}t\right)^n}{n!}\hat{H}^n\Psi(\vec{x},0) &= i\hbar\sum_{n=0}^{\infty}\frac{\partial}{\partial t}\frac{\left(-\frac{i}{\hbar}t\right)^n}{n!}\hat{H}^n\Psi(\vec{x},0) \\
&= i\hbar\sum_{n=0}^{\infty}\frac{n\left(-\frac{i}{\hbar}\right)^nt^{n-1}}{n!}\hat{H}^n\Psi(\vec{x},0) \\
&= \sum_{n=1}^{\infty}\frac{\left(-\frac{i}{\hbar}\right)^nt^{n-1}}{(n-1)!}\hat{H}^n\Psi(\vec{x},0)  \\
&= \sum_{s=0}^{\infty}\frac{\left(-\frac{i}{\hbar}\right)^{s+1}t^s}{s!}\hat{H}^{s+1}\Psi(\vec{x},0) \\
&= \hat{H}\sum_{s=0}^{\infty}\frac{\left(-\frac{i}{\hbar}t\right)^s}{s!}\hat{H}^s\Psi(\vec{x},0) \\
&= \hat{H}\exp\left(-\frac{i}{\hbar}\hat{H}t\right)\Psi(\vec{x},0)
\end{align}

This confirms that $\exp\left(-\frac{i}{\hbar}\hat{H}t\right)\Psi(\vec{x},0)$ satisfies the Schrödinger equation.

For systems with continuous energy spectra (rather than discrete eigenvalues), the superposition becomes an integral:

\begin{equation}
\int_V c(E)\exp\left(-\frac{i}{\hbar}\hat{H}t\right)dE
\end{equation}

This continuous spectrum formulation is particularly relevant for unbound states, such as free particles, where energy can take any value in a continuous range.
\section{Free particle}

The free particle represents a foundational problem in quantum mechanics. While classically a free particle follows a simple trajectory:

\begin{equation}
\vec{r}(t) = \vec{r}_0 + \vec{v}t
\end{equation}

In quantum mechanics, we must solve the Schrödinger equation with the Hamiltonian:

\begin{equation}
\hat{H} = -\frac{\hbar^2}{2m}\nabla^2
\end{equation}

where the potential energy term is zero. The time-independent Schrödinger equation becomes:

\begin{equation}
\hat{H}\Psi = E\Psi
\end{equation}

For a free particle, the energy spectrum is continuous, and the energy depends on the wave vector $\vec{k}$, which is related to momentum by $\vec{p} = \hbar\vec{k}$. The stationary solutions have the form:

\begin{equation}
\Psi_{E(\vec{k})} = e^{-i\frac{E(\vec{k})}{\hbar}t}\Psi_k = e^{-i\frac{E(\vec{k})}{\hbar}t}\frac{e^{i\vec{k}\cdot\vec{x}}}{\sqrt{2\pi}}
\end{equation}

The normalization factor $\frac{1}{\sqrt{2\pi}}$ ensures proper normalization in the context of continuous spectra, where we work with wave packets rather than individual plane waves.

The general solution is formed by superposing these plane wave solutions across all possible wave vectors:

\begin{align}
\Psi &= \int_V d^3k\,a(k)e^{-i\frac{E(\vec{k})}{\hbar}t}\Psi_k \\
&= \int_V d^3k\,a(k)e^{-i\frac{E(\vec{k})}{\hbar}t}\frac{1}{\sqrt{2\pi}}e^{i\vec{k}\cdot\vec{x}}  \\
&= \int_V d^3k\,a(k)\frac{1}{\sqrt{2\pi}}e^{i\left(\vec{k}\cdot\vec{x} - \frac{\hbar\vec{k}^2}{2m}t\right)}
\end{align}

Defining the angular frequency $\omega := \frac{\hbar\vec{k}^2}{2m}$, which represents the energy divided by $\hbar$:

\begin{equation}
\Psi = \int_V d^3k\,a(k)\frac{1}{\sqrt{2\pi}}e^{i(\vec{k}\cdot\vec{x} - \omega t)} \label{eq:3.6}
\end{equation}

This integral represents a superposition of plane waves with different wave vectors, each propagating with its own frequency determined by the dispersion relation $\omega = \frac{\hbar k^2}{2m}$. The coefficient function $a(k)$ determines the contribution of each plane wave to the overall wavefunction.

Unlike classical particles, a quantum free particle is described by a wave packet that generally spreads over time. Even if initially localized, the wave packet will disperse as different momentum components travel at different speeds, a distinctly quantum phenomenon with no classical analog.

The probability density $|\Psi|^2$ evolves in time, showing how the particle's position becomes increasingly uncertain—a manifestation of Heisenberg's uncertainty principle connecting position and momentum uncertainties.


Verifying that our proposed solution satisfies the Schrödinger equation:

For the time derivative:
\begin{align}
i\hbar\frac{\partial\Psi}{\partial t} &= i\hbar\frac{\partial}{\partial t}\int_V d^3k\,a(k)\frac{1}{\sqrt{2\pi}}e^{i(\vec{k}\cdot\vec{x}-\omega t)} \\
&= \int_V d^3k\,a(k)\frac{1}{\sqrt{2\pi}}(-\hbar\omega i)e^{i(\vec{k}\cdot\vec{x}-\omega t)} \\
&= \int_V d^3k\,a(k)\frac{1}{\sqrt{2\pi}}\hbar\omega e^{i(\vec{k}\cdot\vec{x}-\omega t)} \label{eq:3.7} \\
&= \int_V d^3k\,a(k)\frac{1}{\sqrt{2\pi}}\left(\frac{\hbar^2\vec{k}^2}{2m}\right)e^{i(\vec{k}\cdot\vec{x}-\omega t)}
\end{align}

For the Hamiltonian acting on the wavefunction:
\begin{align}
\hat{H}\Psi &= \hat{H}\int_V d^3k\,a(k)\frac{1}{\sqrt{2\pi}}e^{i(\vec{k}\cdot\vec{x}-\omega t)} \\
&= \int_V d^3k\,a(k)\frac{1}{\sqrt{2\pi}}\hat{H}e^{i(\vec{k}\cdot\vec{x}-\omega t)} \\
&= \int_V d^3k\,a(k)\frac{1}{\sqrt{2\pi}}\left(-\frac{\hbar^2}{2m}\nabla^2\right)e^{i(\vec{k}\cdot\vec{x}-\omega t)} \label{eq:3.8} \\
&= \int_V d^3k\,a(k)\frac{1}{\sqrt{2\pi}}\left(\frac{\hbar^2\vec{k}^2}{2m}\right)e^{i(\vec{k}\cdot\vec{x}-\omega t)}
\end{align}

Since the results from equations \eqref{eq:3.7} and \eqref{eq:3.8} are identical, our solution satisfies the Schrödinger equation.

Now we consider the physical meaning of the coefficient function $a(k)$. For an electron beam from an emitter, $a(k)$ represents a distribution of wave vectors centered around a classical momentum $\vec{q}$. This can be modeled as a Gaussian distribution:

\begin{equation}
a(k) = Ae^{-d^2(\vec{k}-\vec{q})^2}
\end{equation}

When $\vec{k} \sim \vec{q}$, the particle behaves approximately like a classical particle. Substituting this expression for $a(k)$ into equation \eqref{eq:3.6}:

\begin{equation}
\Psi = \int_{\mathbb{R}^3}d^3k\,Ae^{-d^2(\vec{k}-\vec{q})^2}\frac{1}{\sqrt{2\pi}}e^{i(\vec{k}\cdot\vec{x}-\omega t)}
\end{equation}

Since the contributions along each dimension are independent, we can factorize this into three one-dimensional integrals:

\begin{equation}
\Psi = \int_{\mathbb{R}^3}\prod_{j=1}^3 dk_j\,A_je^{-d^2(k_j-q_j)^2}\frac{1}{\sqrt{2\pi}}e^{i(k_jx_j-\frac{\hbar k_j^2}{2m}t)}
\end{equation}


To solve this integral, we introduce a simplification by defining:

\begin{equation}
s := \frac{\hbar t}{2m}
\end{equation}

This allows us to rewrite our wavefunction as:

\begin{align}
\Psi &= \int_{\mathbb{R}^3}\prod_{j=1}^3dk_j\,A_je^{-d^2(k_j-q_j)^2}\frac{1}{\sqrt{2\pi}}e^{i(k_jx_j-sk_j^2)} \\
&= \prod_{j=1}^3\frac{A_j}{\sqrt{2\pi}}\underbrace{\int_{\mathbb{R}}dk_j\,e^{-d^2(k_j-q_j)^2+i(k_jx_j-sk_j^2)}}_{I}
\end{align}

To evaluate the integral $I$, we drop the indices temporarily for clarity:

\begin{align}
\phi &= \int_{\mathbb{R}}dk\,e^{-d^2k^2+2d^2kq-d^2q^2+ikx-isk^2} \\
&= e^{-d^2q^2}\int_{\mathbb{R}}dk\,e^{-d^2k^2+2d^2kq+ikx-isk^2}  \\
&= e^{-d^2q^2}\int_{\mathbb{R}}dk\,e^{-k^2(d^2+is)+k(2d^2q+ix)}
\end{align}

This is a Gaussian integral that we can evaluate by defining:

\begin{align}
\nu &:= d^2+is \\
\eta &:= 2d^2q+ix
\end{align}

To complete the square in the exponent, we add and subtract $\frac{\eta^2}{4\nu}$:

\begin{align}
\phi &= e^{-d^2q^2}\int_{\mathbb{R}}dk\,e^{-\nu k^2+\eta k} \\
&= e^{-d^2q^2}\int_{\mathbb{R}}dk\,e^{-\nu k^2+\eta k-\frac{\eta^2}{4\nu}+\frac{\eta^2}{4\nu}}  \\
&= e^{-d^2q^2+\frac{\eta^2}{4\nu}}\int_{\mathbb{R}}dk\,e^{-\nu\left(k-\frac{\eta}{2\nu}\right)^2}
\end{align}

This integral is a standard Gaussian form that can be evaluated using various methods, such as Feynman's trick or polar coordinates in two dimensions.


The solution to the Gaussian integral gives us:

\begin{equation}
\phi = e^{-d^2q^2+\frac{\eta^2}{4\nu}}\sqrt{\frac{\pi}{\nu}}
\end{equation}

Since this expression generally contains complex values, we need to calculate the probability density $\rho$ by evaluating:

\begin{equation}
\rho = |\Psi|^2 = \left|\frac{A_1}{\sqrt{2\pi}}\phi_1\right|^2\left|\frac{A_2}{\sqrt{2\pi}}\phi_2\right|^2\left|\frac{A_3}{\sqrt{2\pi}}\phi_3\right|^2
\end{equation}

For this, we need to compute:

\begin{gather}
|\phi|^2 = \left|\sqrt{\frac{\pi}{\nu}}e^{\text{Re}(z)+i\text{Im}(z)}\right|^2 = \frac{\pi}{|\nu|}e^{2\text{Re}(z)}  \\
-d^2q^2+\frac{\eta^2}{4\nu} = -d^2q^2+\frac{(d^2-is)(2d^2q+ix)^2}{4(d^4+s^2)} = \\
= -d^2q^2+\frac{(d^2-is)(4d^4q^2+4id^2qx-x^2)}{4(d^4+s^2)} = \\
= -d^2q^2+\frac{d^2(4d^4q^2+4id^2qx-x^2)-is(4d^4q^2+4id^2qx-x^2)}{4(d^4+s^2)} =  \\
= -d^2q^2+\frac{d^2(4d^4q^2+4sqx-x^2)}{4(d^4+s^2)}-i\frac{(4sd^4q^2+4d^4qx-sx^2)}{4(d^4+s^2)}
\end{gather}

Therefore:

\begin{align}
\text{Re}(z) &= -d^2q^2+\frac{d^2(4d^4q^2+4sqx-x^2)}{4(d^4+s^2)} \\
&= -d^2q^2+\frac{d^2(4d^4q^2+4s^2q^2-4s^2q^2+4sqx-x^2)}{4(d^4+s^2)} \\
&= -d^2q^2+\frac{d^2(4d^4q^2+4s^2q^2-(x-2sq)^2)}{4(d^4+s^2)} \\
&= -d^2q^2+\frac{4d^2q^2(d^4+s^2)-d^2(x-2sq)^2}{4(d^4+s^2)}  \\
&= -d^2q^2+d^2q^2-\frac{d^2(x-2sq)^2}{4(d^4+s^2)} \\
&= -\frac{d^2(x-2sq)^2}{4(d^4+s^2)}
\end{align}

Substituting the value of $s$ back:

\begin{equation}
2sq = 2\frac{\hbar t}{2m}q = \frac{\hbar q}{m}t = vt
\end{equation}

where $v = \frac{\hbar q}{m}$ is the velocity.

Finally, we can evaluate one component of the probability density:

\begin{align}
\left|\frac{A_j}{\sqrt{2\pi}}\phi_j\right|^2 &= \frac{|A_j|^2}{2\pi}\frac{\pi}{|\nu|}e^{-\frac{d^2(x_j-v_jt)^2}{4(d^4+s^2)}} \\
&= \frac{|A_j|^2}{2\sqrt{d^4+s^2}}e^{-\frac{d^2(x_j-v_jt)^2}{4(d^4+s^2)}}
\end{align}

This represents a Gaussian wave packet that moves with velocity $v_j$ along the $x_j$-axis while spreading out over time. At $t=0$, it is centered at $x=0$, and as time progresses, it both translates and spreads until it becomes essentially flat.
